\chapter{Installazione e Avvio}

\section{Prerequisiti}

\begin{itemize}
  \item \textbf{Docker} $\geq$ 24.0 e \textbf{Docker Compose} $\geq$ 2.20
  \item \textbf{make} (opzionale, semplifica i comandi)
  \item \textbf{Chiave API OpenAI} (se \texttt{LLM\_PROVIDER=openai}) oppure
        \textbf{Ollama} installato localmente (se \texttt{LLM\_PROVIDER=ollama})
  \item Connessione Internet per il pull delle immagini Docker al primo avvio
\end{itemize}

\section{Struttura del Repository}

\begin{lstlisting}[language=bash]
ai-document-assistant/
+-- app/
|   +-- main.py          # Entry point FastAPI
|   +-- config.py        # Impostazioni (pydantic-settings)
|   +-- api/             # Router: documents, chat, search
|   +-- core/            # Pipeline: ingestion, retrieval, llm
|   +-- models/          # Modelli SQLAlchemy
|   +-- db/              # Session, CRUD
|   +-- utils/           # Validazione file
+-- docker/
|   +-- Dockerfile
|   +-- docker-compose.yml
|   +-- .env.example
+-- docs/                # Questa documentazione (LaTeX)
+-- tests/
|   +-- fixtures/sample.pdf
|   +-- test_api.py
|   +-- test_ingestion.py
+-- requirements.txt
+-- Makefile
\end{lstlisting}

\section{Configurazione}

\begin{enumerate}
  \item Copiare il file di esempio:
\begin{lstlisting}[language=bash]
cp docker/.env.example docker/.env
\end{lstlisting}

  \item Modificare \texttt{docker/.env} con i propri valori:
\end{enumerate}

\begin{lstlisting}[language=bash]
# Scegliere il provider LLM
LLM_PROVIDER=openai          # oppure: ollama

# Obbligatorio se LLM_PROVIDER=openai
OPENAI_API_KEY=sk-...

# Dimensione vettori: 1536 (OpenAI) oppure 768 (nomic-embed-text)
QDRANT_VECTOR_SIZE=1536

# Password PostgreSQL (cambiare in produzione)
POSTGRES_PASSWORD=password
\end{lstlisting}

\textbf{Nota per Ollama:} decommentare il servizio \texttt{ollama} nel
\texttt{docker-compose.yml} e impostare \texttt{LLM\_PROVIDER=ollama}.
Al primo avvio eseguire il pull del modello:
\begin{lstlisting}[language=bash]
docker exec ai_doc_ollama ollama pull llama3.2
docker exec ai_doc_ollama ollama pull nomic-embed-text
\end{lstlisting}

\section{Avvio con Docker Compose}

\begin{lstlisting}[language=bash]
# Build e avvio di tutti i servizi
docker compose -f docker/docker-compose.yml up --build -d

# Oppure, usando il Makefile:
make build
make up
\end{lstlisting}

I servizi avviati sono:

\begin{table}[H]
  \centering
  \begin{tabularx}{\textwidth}{llX}
    \toprule
    \textbf{Servizio} & \textbf{Porta} & \textbf{URL} \\
    \midrule
    FastAPI app & 8000 & \url{http://localhost:8000} \\
    Swagger UI  & 8000 & \url{http://localhost:8000/docs} \\
    PostgreSQL  & 5432 & \texttt{localhost:5432} \\
    Qdrant REST & 6333 & \url{http://localhost:6333} \\
    Qdrant gRPC & 6334 & \texttt{localhost:6334} \\
    \bottomrule
  \end{tabularx}
  \caption{Porte esposte dai servizi}
\end{table}

\section{Verifica dell'Installazione}

\begin{lstlisting}[language=bash]
# Health check
curl http://localhost:8000/health
# {"status": "ok", "version": "1.0.0"}

# Carica un PDF di test
curl -X POST http://localhost:8000/documents/upload \
  -F "file=@tests/fixtures/sample.pdf"

# Fai una domanda
curl -X POST http://localhost:8000/chat/ \
  -H "Content-Type: application/json" \
  -d '{"question": "Di cosa tratta il documento?"}'
\end{lstlisting}

\section{Comandi Make}

\begin{table}[H]
  \centering
  \begin{tabularx}{\textwidth}{lX}
    \toprule
    \textbf{Comando} & \textbf{Effetto} \\
    \midrule
    \texttt{make up}      & Avvia i container in background \\
    \texttt{make down}    & Ferma i container \\
    \texttt{make build}   & Ricostruisce le immagini \\
    \texttt{make logs}    & Mostra i log del servizio app \\
    \texttt{make test}    & Esegue pytest all'interno del container \\
    \texttt{make shell}   & Apre una shell bash nel container app \\
    \texttt{make clean}   & Rimuove container, volumi e cache Python \\
    \texttt{make docs}    & Compila la documentazione LaTeX in PDF \\
    \bottomrule
  \end{tabularx}
  \caption{Comandi disponibili nel Makefile}
\end{table}

\section{Variabili d'Ambiente}

\begin{longtable}{l p{3.2cm} p{7cm}}
  \toprule
  \textbf{Variabile} & \textbf{Default} & \textbf{Descrizione} \\
  \midrule
  \endfirsthead
  \midrule
  \textbf{Variabile} & \textbf{Default} & \textbf{Descrizione} \\
  \midrule
  \endhead
  \bottomrule
  \endfoot
  \bottomrule
  \caption{Variabili d'ambiente configurabili}
  \endlastfoot
  \texttt{LLM\_PROVIDER}           & \texttt{openai}             & Provider LLM: \texttt{openai} o \texttt{ollama} \\
  \texttt{OPENAI\_API\_KEY}         & —                           & Chiave API OpenAI (obbligatoria se openai) \\
  \texttt{OPENAI\_MODEL}            & \texttt{gpt-4o-mini}        & Modello chat OpenAI \\
  \texttt{OPENAI\_EMBEDDING\_MODEL} & \texttt{text-embedding-3-small} & Modello embedding OpenAI \\
  \texttt{OLLAMA\_BASE\_URL}        & \texttt{http://ollama:11434} & URL servizio Ollama \\
  \texttt{OLLAMA\_MODEL}            & \texttt{llama3.2}           & Modello chat Ollama \\
  \texttt{OLLAMA\_EMBEDDING\_MODEL} & \texttt{nomic-embed-text}   & Modello embedding Ollama \\
  \texttt{QDRANT\_URL}              & \texttt{http://qdrant:6333} & URL servizio Qdrant \\
  \texttt{QDRANT\_COLLECTION}       & \texttt{doc\_chunks}        & Nome collection Qdrant \\
  \texttt{QDRANT\_VECTOR\_SIZE}     & \texttt{1536}               & Dimensione vettori embedding \\
  \texttt{DATABASE\_URL}            & —                           & URL asyncpg PostgreSQL \\
  \texttt{CHUNK\_SIZE}              & \texttt{500}                & Dimensione chunk in caratteri \\
  \texttt{CHUNK\_OVERLAP}           & \texttt{50}                 & Overlap tra chunk consecutivi \\
  \texttt{TOP\_K}                   & \texttt{4}                  & Numero chunk recuperati per query \\
  \texttt{MAX\_UPLOAD\_SIZE\_MB}    & \texttt{50}                 & Limite dimensione file PDF \\
  \texttt{DEBUG}                    & \texttt{false}              & Modalità debug (SQL logging) \\
\end{longtable}

\section{Conclusioni e Sviluppi Futuri}

Il progetto raggiunge tutti gli obiettivi tecnici prefissati: un sistema RAG funzionante,
containerizzato, con API documentata e supporto per provider LLM multipli.

\subsection{Possibili Miglioramenti}

\begin{itemize}
  \item \textbf{Autenticazione}: aggiungere JWT/OAuth2 per proteggere gli endpoint.
  \item \textbf{Re-ranking}: aggiungere un nodo LangGraph di re-ranking (es.\ Cohere
        Rerank) per migliorare la qualità dei chunk recuperati.
  \item \textbf{Interfaccia Web}: sviluppare un frontend React/Next.js con chat UI.
  \item \textbf{Supporto multi-formato}: estendere l'ingestion a DOCX, PPTX, HTML.
  \item \textbf{Chunking semantico}: usare semantic chunking basato su boundary
        semantici invece di caratteri.
  \item \textbf{Cache embedding}: implementare una cache Redis per evitare il ricalcolo
        di embedding identici.
  \item \textbf{Monitoring}: integrare Prometheus/Grafana per metriche di latenza e
        throughput.
  \item \textbf{Multi-tenant}: supporto per utenti multipli con isolamento dei documenti.
\end{itemize}
